\section{Introduction}
Three-dimensional and multi-die FPGAs expand the design space of programmable interconnects, while placing new demands on the CAD tools used to evaluate them. Recent frameworks support architecture generation, 3D placement and multi-die routing~\cite{lazagna,flatland,multidie}. Lookahead tables summarize delay/congestion costs by resource and destination context. Construction matters when architecture or cost changes invalidate an existing table. Compatible precomputed tables should be reused; our focus is the construction work that remains when such reuse is unavailable.

However, parallel routing does not automatically make lookahead construction parallel. In our locked VPR flow, an eight-core GEMM experiment spends 149.56 seconds constructing lookahead---approximately 45\% of its fixed-input process time. Two execution bottlenecks restrict concurrency: barriers between groups containing only a few samples, and fine-grained priority ordering within each search. Simply merging searches or relaxing their ordering is not a drop-in solution. Source/layer-relative recording context and associated delay/congestion representatives must survive reordering. The challenge is to separate removable scheduling constraints from output-defining dependencies.

MultiQueue-based routing parallelizes connection searches through carefully controlled priority relaxation~\cite{multiqueue}; it does not directly provide the sample-recording and reduction semantics of lookahead construction. Multi-die-aware lookahead improves how heuristic costs represent die crossings~\cite{multidie}, addressing the heuristic itself rather than the execution organization studied here. Parallel shortest-path methods already exploit buckets and frontier expansion to balance useful work against scheduling overhead~\cite{delta,gpu_sssp}. Lookahead additionally requires source-relative cells, coupled labels, same-bucket completion and native integration. Caching and sparse sampling can avoid or reduce construction~\cite{vproptions}; we parallelize the retained searches.

Our insight separates logical sample identity from execution order. STRICT schedules private samples across groups and reduces in native order under explicit equivalence conditions. ALT separately defines FP32 whole-pair waves within searches. Both integrate through logical sample/context boundaries, with CPU scheduling across samples and GPU waves within samples. The measured benefits extend beyond isolated kernels: STRICT Pipeline achieves a $1.638\times$ Process speedup over Stock on the eight-core GEMM cohort, while ALT CUDA achieves $3.207\times$ native construction and $1.094\times$ Process speedups over optimized CPU Pipeline on the same-family LSTM tuning holdout. These results are from separate cohorts. Our contributions are (C1) source-context-preserving cross-group decoupling with explicit reduction conditions, (C2) whole-pair FP32 wave execution with explicit closure and numeric limits, and (C3) native CPU/GPU integration with mechanism controls, route quality validation and measured applicability limits.
